An object of this class represents a program module, which is part of
a program's executable image.
A \code{BPatch\_module} represents a source file in either an
executable or a shared library. Dyninst automatically creates
a module called \code{DEFAULT\_MODULE} in each executable to hold any
objects that it cannot match to a source file.
\code{BPatch\_module} objects are obtained by calling the
\code{BPatch\_image} member function \code{getModules}.

\begin{apient}
  std::vector<BPatch_function*> *findFunction(const char* name,
                                              std::vector<BPatch_function*> &funcs,
                                              bool notify_on_failure = true,
                                              bool regex_case_sensitive = true,
                                              bool incUninstrumentable = false,
                                              bool dont_use_regex = false);
\end{apient}

\apidesc{
  Return a vector of \code{BPatch\_functions} matching name, or NULL if the function does not
  exist. If name contains a POSIX-extended regular expression, a regex search will be per-
  formed on function names, and matching \code{BPatch\_functions} returned. If the \code{incUninstrumentable}
  flag is set, the returned table of procedures will include
  uninstrumentable functions. The default behavior is to omit these functions.
  [NOTE: If name is not found to match any demangled function names in the module, the
  search is repeated as if name is a mangled function name. If this second search succeeds,
  functions with mangled names matching name are returned instead.]
}

\begin{apient}
  BPatch_Vector<BPatch_function *> *findFunctionByAddress(void* addr,
                                                          BPatch_Vector<BPatch_function *> &funcs,
                                                          bool notify_on_failure = true,
                                                          bool incUninstrumentable = false);
\end{apient}

\apidesc{
  Return a vector of \code{BPatch\_function} s that contains addr, or
  NULL if the function does not exist.
  If the incUninstrumentable flag is set, the returned table of
  procedures will include uninstrumentable functions. The
  default behavior is to omit these functions.
}

\begin{apient}
  BPatch_function* findFunctionByEntry(Dyninst::Address addr);
\end{apient}

\apidesc{
  Returns the function that begins at the specified address addr.
}

\begin{apient}
  BPatch_function* findFunctionByMangled(const char* mangled_name,
                                         bool incUninstrumentable = false);
\end{apient}

\apidesc{
  Return a \code{BPatch\_function} for the mangled function name
  defined in the module corresponding to the invoking
  \code{BPatch\_module}, or NULL if it does not define the function.

  If the incUninstrumentable flag is set, the functions searched will
  include uninstrumentable functions. The default
  behavior is to omit these functions.
}

\begin{apient}
  bool getAddressRanges(const char* fileName, unsigned int lineNo,
                        std::vector<std::pair<unsigned long, unsigned long>> &ranges);
\end{apient}

\apidesc{
  Given a filename and line number, fileName and lineNo, this
  function returns the ranges of mutatee
  addresses that implement the code range in the output parameter
  ranges. In many cases a source code line will only
  have one address range implementing it. However, compiler
  optimizations may turn this into multiple, disjoint address
  ranges. This information is only available if the mutatee was
  compiled with debug information.
  This function may be more efficient than the \code{BPatch\_process}
  version of this function. Calling
  \code{BPatch\_process::getAddressRange} will cause Dyninst to parse
  line information for all modules in a process. If
  \code{BPatch\_module::getAddressRange} is called then only the
  debug information in this module will be parsed.
  This function returns true if it was able to find any line
  information, false otherwise.
}

\begin{apient}
  size_t getAddressWidth();
\end{apient}

\apidesc{
  Return the size (in bytes) of a pointer in this module. On 32-bit
  systems this function will return 4, and on 64-bit
  systems this function will return 8.
}

\begin{apient}
  void* getBaseAddr();
\end{apient}

\apidesc{
  Return the base address of the module. This address is defined as
  the start of the first function in the module.
}

\begin{apient}
  std::vector<BPatch_function *> *getProcedures(bool incUninstrumentable = false);
\end{apient}

\apidesc{
  Return a vector containing the functions in the module.
}

\begin{apient}
  char* getFullName(char* buffer, int length);
\end{apient}

\apidesc{
  Fills buffer with the full path name of a module, up to length
  characters when this information is available.
}

\begin{apient}
  BPatch_hybridMode getHybridMode();
\end{apient}

\apidesc{
  Return the mutator\'s analysis mode for the mutate; the default
  mode is the normal mode.
}

\begin{apient}
  char* getName(char* buffer, int len);
\end{apient}

\apidesc{
  This function copies the filename of the module into buffer, up to
  len characters. It returns the value of the buffer
  parameter.
}

\begin{apient}
  unsigned long getSize();
\end{apient}

\apidesc{
  Return the size of the module. The size is defined as the end of
  the last function minus the start of the first
  function.
}

\begin{apient}
  bool getSourceLines(unsigned long addr, std::vector<BPatch_statement> &lines);
\end{apient}

\apidesc{
  This function returns the line information associated with the
  mutatee address addr. The vector lines contain pairs of
  filenames and line numbers that are associated with addr. In many
  cases only one filename and line number is
  associated with an address, but certain compiler optimizations may
  lead to multiple filenames and lines at an address.
  This information is only available if the mutatee was compiled with
  debug information.
  This function may be more efficient than the \code{BPatch\_process}
  version of this function. Calling
  \code{BPatch\_process::getSourceLines} will cause Dyninst to parse
  line information for all modules in a process. If
  \code{BPatch\_module::getSourceLines} is called then only the debug
  information in this module will be parsed.
  This function returns true if it was able to find any line
  information at addr, or false otherwise.
}

\begin{apient}
  bool getVariables(std::vector<BPatch_variableExpr *> &vars);
\end{apient}

\apidesc{
  Fill the vector vars with the global variables that are specified
  in this module. Returns false if no results are
  found and true otherwise.
}

\begin{apient}
  bool isExploratoryModeOn();
\end{apient}

\apidesc{
  This function returns true if the mutator's analysis mode sets to
  the defensive mode or the
  exploratory mode.
}

\begin{apient}
  bool isSharedLib();
\end{apient}

\apidesc{
  This function returns true if the module is part of a shared library.
}
